\section[\texorpdfstring{\(L^p\)}{Lp} spaces]{\texorpdfstring{\(L^p\)}{Lp} spaces  \marginpar{\red{Lecture 2}}}
\label{sec: Lp spaces}

\begin{definition}
    We define the \(\cL^p\) space to be
    \[
        \cL^p \coloneq \cL^p(\domain) \coloneq \set{X\colon \Omega \to \domain : \E{\norm{X}^p} < \infty}
    \]
    and define the seminorm for \(X\in \cL^p\) as
    \[
        \norm{X}_{L^p} \coloneq \E[\big]{\norm{X}^p}^{\frac1p}.
    \]
\end{definition}

We called \(\norm{\cdot}_{L^p}\) a ``seminorm''. For this we need to prove that
it actually satisfies the seminorm properties. This is the content of the
following theorem.

\begin{theorem}[Seminorm]
    \label{thm: seminorm properties}
    \(\norm{\cdot}_{L^p}\) satisfies the seminorm properties. That is for \(X,Y
    \in \cL^p(\domain)\) with \((\domain, \norm{\cdot})\) a normed space we have
    \begin{enumerate}[label=(\roman*)]
        \item\label{it: positive}
        \(\norm{X}_{L^p} \ge 0\),
        \hfill (positive)
        \item\label{it: scaling}
        \(\norm{\lambda X}_{L^p} = |\lambda| \norm{X}_{L^p}\)
        for all \(\lambda \in \real\),
        \hfill (scaling)
        \item\label{it: triangle inequality}
        \(\norm{X + Y}_{L^p} \le \norm{X}_{L^p} + \norm{Y}_{L^p}\).
        \hfill (triangle/Minkowski inequality)
    \end{enumerate}
    Moreover we have for any \(p \le q\)
    \begin{equation}
        \label{eq: Lp norm domination}
        \norm{X}_{L^p} \le \norm{X}_{L^q}.
    \end{equation}
\end{theorem}

\begin{remark}[\(L^p\) convergence]
    \label{rem: Lp convergence}
    Observe that \(X_n \xrightarrow{L^p} X\) (Definition \ref{def: types of convergence}) if and only if \(\norm{X_n -X}_{L^p} \to
    0\). And Equation \eqref{eq: Lp norm domination} means for \(1\le p\le q\) that \(\norm{X_n - X}_{L^q} \to 0\) implies \(\norm{X_n - X}_{L^p} \to 0\).
\end{remark}

We will prove the semi-norm properties later as a special case of the triangle
inequality for Orlicz norms (Theorem \ref{thm: orlicz seminorm}). For now we
will only prove the norm domination property \eqref{eq: Lp norm domination}.
For this we require Hölder's inequality.

\begin{theorem}[Hölder's inequality]
    \label{thm: holder's inequality}
    For \(r,s >1\) with \(\frac1s + \frac1r = 1\) we have
    \[
        \E{\abs{XY}} \le \E{\abs{X}^r}^{\frac1r}\E{\abs{Y}^s}^{\frac1s} = \norm{X}_{L^r} \norm{Y}_{L^s}.
    \]
    This also holds true if the quantities are not finite!
\end{theorem}

You will prove this inequality in Exercise \ref{ex: hölder's inequality}.

\begin{proof}[Proof of \eqref{eq: Lp norm domination}]
    Assume without loss of generality the strict
    inequality \(p< q\) and define \(r
    \coloneq \frac{q}p>1\) and \(s \coloneq \frac{p}{q-p}>1\) such
    that \(\frac1{r} + \frac1{s} = 1\). Then
    \[
    \begin{aligned}
        \norm{X}_{L^p}^p = \E{\norm{X}^p}
        &= \E{\abs{\norm{X}^p \cdot 1}}
        \\
        \overset{\text{Hölder}}&\le \E{\norm{X}^{pr}}^{\frac1r} \E{1^s}^{\frac1s}
        = \E{\norm{X}^q}^{\frac{p}q}
        = \norm{X}_{L^q}^p.
    \end{aligned}\]
    Taking the \(p\)-th root gives \eqref{eq: Lp norm domination}.
\end{proof}

\begin{example}[Not positive definite]
    Observe that the missing property for \(\norm{\cdot}_{L^p}\) to be a norm 
    is that \(\norm{X}_{L^p} = 0\) implies \(X=0\). As an example consider
    a uniformly distributed \(U\sim \uniform(0,1)\) with density \(f(x) = \ind{[0,1]}\).
    Define
    \[
        X \coloneq \begin{cases}
            U & U \neq 1
            \\
            0 & U = 1.
        \end{cases}
    \]
    Then clearly \(U\neq X\) but for any \(p\ge 1\)
    \[
        \E{\abs{X-U}^p}
        = \int_0^1 \abs{x^p \ind{[0,1)}(x) - x^p} dx
        = \int_0^1 \ind{\set{1}} dx = 0,
    \]
    and consequently \(\norm{X-U}_{L^p} = 0\).
\end{example}

\begin{fact}
    \label{fact: Lp zero norm}
    \(\norm{X-Y}_{L^p} = 0\) if and only if \(\prob{X=Y}=1 \overset{\text{def.}}\iff X\overset{\as}= Y\).
\end{fact}
\begin{proof}[Proofsketch]
    Measure theory: \(f\ge 0\) and \(\int f d\mu = 0\) 
    implies \(\mu(f\neq 0) = 0\).
\end{proof}

Using this inequality we can now prove Theorem \ref{thm: seminorm properties}


\begin{definition}[\(L^p\) space]
    We define the normed space \((L^p,\norm{\cdot}_{L^p})\) as
    \[
        L^p \coloneq L^p(\domain) \coloneq \set{[X]: X \in \cL^p(\domain)},
        \quad\text{with}\quad
        \norm{[X]}_{L^p} \coloneq \norm{X}_{L^p}
    \]
    where \([X]\) is the equivalence class of the equivalence relation
    \[
        X\sim Y :\iff \norm{X-Y}_{L^p} = 0 \iff X\overset{\as}= Y.
    \]
\end{definition}
\begin{remark}[Norm well defined]
    The definition of the norm is well-defined as for \(X,Y\in [X]\)
    we have by the triangle inequality
    \[
        \norm{X}_{L^p} \le \norm{Y}_{L^p} + \norm{X - Y}_{L^p} \overset{X,Y\in [X]}= \norm{Y}_{L^p}
        \overset{\text{similarly}}\le \norm{X}_{L^p}.
    \]
\end{remark}

It turns out that proving the seminorm properties (Theorem \ref{thm: seminorm
properties}), is easier in a more general setting. Specifically, since
deterministic values may be moved outside of the expectation, we have
\[
    \norm{X}_{L^p} = \inf\set[\Big]{\lambda > 0 : \E[\Big]{\bigl(\tfrac{\norm{X}}{\lambda}\bigr)^p} \le 1}.
\]
With \(\psi(x) = x^p\) the \(L^p\) norm is thus a special case of
an ``Orlicz norm''.
\begin{definition}[Orlicz norm]
    \label{def: orlicz norm}
    Let \(\psi\colon [0,\infty) \to [0,\infty)\) be a non-decreasing, convex function
    with \(\psi(0) = 0\) and \(\lim_{x\to \infty} \psi(x) = \infty\). Then we define
    for \(X\colon \Omega \to \domain\) the \textbf{Orlicz norm}
    \[
        \norm{X}_{\psi} \coloneq \inf\set[\Big]{\lambda > 0 : \E[\Big]{\psi\bigl(\tfrac{\norm{X}}{\lambda}\bigr)} \le 1}.
    \]
    The \textbf{Orlicz spaces} are given by \(\cL^\psi \coloneq \set{X\colon \Omega \to \domain : \norm{X}_{\psi} < \infty}\).
\end{definition}

\begin{theorem}
    \label{thm: orlicz seminorm}
    The Orlicz norm is a seminorm.
\end{theorem}
\begin{proof}
    The Orlicz norm is non-negative by definition of the infimum. The scaling
    property follows from
    \[\begin{aligned}
        \norm{\alpha X}_{\psi}
        &= \inf\set[\Big]{\lambda > 0 : \E[\Big]{\psi\bigl(\tfrac{\abs{\alpha}\norm{X}}{\lambda}\bigr)} \le 1}
        \\
        \overset{\mu=\frac{\lambda}{\abs{\alpha}}}&= \abs{\alpha} \inf\set[\Big]{\mu > 0 : \E[\Big]{\psi\bigl(\tfrac{\norm{X}}{\mu}\bigr)} \le 1}
        &= |\alpha| \norm{X}_{\psi}.
    \end{aligned}
    \]

    Let \(a\coloneq \norm{X}_{\psi}\) and \(b\coloneq \norm{Y}_{\psi}\). Then
    for any \(\epsilon >0\) we have by definition of the Orlicz norm
    \begin{equation}
        \label{eq: orlicz epsilon}     
        \E{\psi\bigl(\tfrac{\norm{X}}{a + \epsilon}\bigr)} \le 1,
        \quad\text{and}\quad
        \E{\psi\bigl(\tfrac{\norm{Y}}{b + \epsilon}\bigr)} \le 1.
    \end{equation}
    Using \(\lambda = \frac{a+\epsilon}{a+b+2\epsilon}\in[0,1]\) and thereby \(1-\lambda = \frac{b+\epsilon}{a+b+2\epsilon}\) we get
    \[
        \frac{\norm{X}+\norm{Y}}{a + b + 2\epsilon} = \lambda \frac{\norm{X}}{a + \epsilon} + (1-\lambda) \frac{\norm{Y}}{b + \epsilon}.
    \]
    Finally, using that \(\psi\) is non-decreasing and convex:
    \[\begin{aligned}
        \E[\Big]{\psi\Bigl(\frac{\norm{X+Y}}{a + b + 2\epsilon}\Bigr)}
        \overset{\text{non-dec.}}&\le \E[\Big]{\psi\Bigl(\frac{\norm{X}+\norm{Y}}{a + b + 2\epsilon}\Bigr)}
        \\
        \overset{\text{conv.}}&\le \lambda \E[\Big]{\psi\Bigl(\frac{X}{a + \epsilon}\Bigr)} + (1-\lambda) \E[\Big]{\psi\Bigl(\frac{Y}{b + \epsilon}\Bigr)}
        \\
        \overset{\ref{eq: orlicz epsilon}}&\le \lambda + (1-\lambda) = 1.
    \end{aligned}
    \]
    Thus by definition of the Orlicz norm we have shown
    \[
        \norm{X + Y}_{\psi} \le a + b + 2\epsilon
        \overset{\text{def.}}= \norm{X}_{\psi} + \norm{Y}_{\psi} + 2\epsilon.
    \]
    As this is true for any \(\epsilon > 0\), we have shown the triangle inequality.
\end{proof}

Again we have that the norm of differences is only zero for almost surely equal
random variables.

\begin{fact}
    \(\norm{X}_\psi = 0\) if and only if \(X\overset{\as}= 0\).
\end{fact}
\begin{proof}[Proofsketch]
    If \(X \overset{\as}= 0\), then \(\psi(0)=0\) implies the Orlicz
    norm is zero. The other direction uses the property \(\lim_{x\to\infty}
    \psi(x) = \infty\) as \(\lambda \to 0\) explodes the input if \(\norm{X}\)
    is non-zero.
\end{proof}

\begin{fact}[\(L^p\) is a Banach space]
    \label{fact: Lp complete}
    If \((\domain, \norm{\cdot})\) is a Banach space, then for every \(p \ge 1\)
    the normed space \((L^p(\domain), \norm{\cdot}_{L^p})\) is
    also a Banach space. That is it is \underline{complete},
    which means that every Cauchy sequence converges in \(L^p\). 
\end{fact}
\begin{proof}
See e.g. \citep[Thm.~6.25 and 7.3]{klenkeProbabilityTheoryComprehensive2014}.
\end{proof}

Note that for any \(X,Y\in [X] \in L^p\) we have \(X\overset{\as}= Y\)
and recall that most limits in probability were unique in the same sense
(Exercise \ref{ex: unique stochastic limits}).
Since most equations in probability theory only hold almost surely, this caveat
is often omitted and one simply writes \(X=Y\), blurring the lines
between the equivalence classes \(L^p\) and the representatives in \(\cL^p\).

\subsection*{Exercises}

\begin{exercise}[Proof of Hölder's inequality \Coffeecup]
    \label{ex: hölder's inequality}
    Let \(\frac1p + \frac1q = 1\) with \(p,q > 1\).
    \begin{enumerate}[label=(\alph*)]
        \item Prove ``Young's inequality'': for \(a,b \ge 0\) we have
        \[
            ab \le \tfrac{a^p}p + \tfrac{b^q}q.
        \]
        \textbf{Hint:} \emph{Option 1}: Let \(x = p \log(a)\) and \(y = q \log(b)\)
        and use convexity of the exponential function.
        \emph{Option 2 (brute force)}: Find the minimum of 
        the function \(f(x) = \frac{x^p}p + \frac{b^q}q - xb\).
        \begin{solution}
            \emph{Option 1:} Using the convexity of the exponential function,
            \[\begin{aligned}
                ab
                = e^{\log(a) + \log(b)}
                = e^{\frac1p (p\log(a)) + \frac1q (q\log(b))}
                &\le \frac1p e^{p\log(a)} + \frac1q e^{q\log(b)}
                \\
                &= \frac{a^p}p + \frac{b^q}q.
            \end{aligned}
            \]
            \emph{Option 2:} On \([0, \infty)\), \(f\) is convex  as \(f''(x) = (p-1)x^{p-2} \ge 0\).
            The minimum is therefore attained at \(f'(x) = x^{p-1} - b = 0\), that is
            at \(x_* = b^{\frac1{p-1}}\). Since \(q= \frac{p}{p-1}\) we have that \(x_*^p = b^q\).
            \[
                f(x_*) = \frac{b^q}p + \frac{b^q}q - b^{\frac{1}{p-1} + 1} = 0.
            \]
            Consequently, for all \(x \ge 0\) we have \(f(x) \ge f(x_*) = 0\).
            In particular for \(x = a\) we have the claim.
        \end{solution}

        \item Prove ``Hölder's inequality'' 
        For \(p,q >1\) with \(\frac1p + \frac1q = 1\) we have
        \[
            \E{\abs{XY}} \le \E{\abs{X}^p}^{\frac1p}\E{\abs{Y}^q}^{\frac1q} = \norm{X}_{L^p} \norm{Y}_{L^q}.
        \]

        \textbf{Hint:} 
        Assume without loss of generality \(X \in \cL^p(\real)\) and \(Y\in \cL^q(\real)\),
        otherwise the right hand side is infinite and the claim is trivial.
        Apply Young's inequality to \(a =
        \frac{\abs{X}}{\norm{X}_{L^p}}\) and \(b = \frac{\abs{Y}}{\norm{Y}_{L^q}}\).
        \begin{solution}
            Applying Young's inequality we have
            \[
                \frac{\abs{X}}{\norm{X}_{L^p}} \cdot \frac{\abs{Y}}{\norm{Y}_{L^q}}
                \le \frac1p \Bigl(\frac{\abs{X}}{\norm{X}_{L^p}}\Bigr)^p
                + \frac1q \Bigl(\frac{\abs{Y}}{\norm{Y}_{L^q}}\Bigr)^q.
            \]
            Taking expectation on both sides we get
            \[
                \frac{\E[\big]{\abs{XY}}}{\norm{X}_{L^p} \norm{Y}_{L^q}}
                \le 
                    \frac1p \E[\Big]{\Bigl(\frac{\abs{X}}{\norm{X}_{L^p}}\Bigr)^p}
                    + \frac1q \E[\Big]{\Bigl(\frac{\abs{Y}}{\norm{Y}_{L^q}}\Bigr)^q}
                = 
                    \frac1p + \frac1q
                = 1,
            \]
            where we use that \(\norm{X}_{L^p}^p = \E{\abs{X}^p}\) is
            deterministic to move it outside of the expectation. Multiplying
            both sides by \(\norm{X}_{L^p} \norm{Y}_{L^q}\)
            implies the claim.
        \end{solution}
    \end{enumerate}
\end{exercise}

\begin{exercise}[\(L^2\)-Hilbertspace \Coffeecup]
    \label{ex: L2 hilbert space}
    Let \(X,Y\) be \(\cL^2\) random variables in a Hilbert space \((\hilbert, \scp{\cdot, \cdot})\).
    Show that 
    \[
        \langle [X], [Y] \rangle_{L^2} \coloneq \langle X, Y\rangle_{L^2} \coloneq \E{\scp{X,Y}}
    \]
    defines an inner product on \(L^2\) which induces the norm \(\norm{\cdot}_{L^2}\).
    \begin{solution}
        Clearly \(\scp{X, X}_{L^2} = \norm{X}_{L^2}^2\) so the inner product induces
        the norm \(\norm{\cdot}_{L^2}\). We need to show the inner product properties:
        \begin{enumerate}[label=(\roman*)]
            \item\label{it: inner product positive}
            \(\scp{X, X}_{L^2} \ge 0\) with equality if and only if \([X]=0 \overset{\text{Fact } \ref{fact: Lp zero norm}}\iff X=0\) a.s.
            \\ \null \hfill (positive definiteness)
            \item\label{it: inner product symmetry}
            \(\scp{X, Y}_{L^2} = \scp{Y, X}_{L^2}\).
            \hfill (symmetry)
            \item\label{it: inner product linearity}
            \(\scp{aX + bY, Z}_{L^2} = a \scp{X, Z}_{L^2} + b \scp{Y, Z}_{L^2}\)
            for all \(a,b\in \real\).
            \hfill (linearity in first argument)
        \end{enumerate}
        Property \ref{it: inner product positive} follows from \(\scp{X, X}_{L^2} = \norm{X}_{L^2}^2\)
        and this property is established in Theorem \ref{thm: seminorm properties} and
        Fact \ref{fact: Lp zero norm}.
        Property \ref{it: inner product symmetry} follows from the
        commutativity of multiplication. Property \ref{it: inner product linearity}
        follows from the linearity of expectation.
        \(L^2\) is further a complete normed space by Fact \ref{fact: Lp complete}.
        Consequently, \((L^2, \scp{\cdot, \cdot}_{L^2})\) is a Hilbert space.
    \end{solution}
\end{exercise}

\begin{exercise}[Expectation -- Projection in \(L^2\) \Coffeecup\Coffeecup]
    \label{ex: projection in L2}
    Recall that by Exercise \ref{ex: L2 hilbert space}, \(L^2\) is a Hilbertspace.
    Identify deterministic random variables with real numbers  such that \(\real
    \subseteq \cL^2(\real)\). Then identify \(\real\) with \(\set{[x]:
    x\in\real} \subseteq L^2\) such that \(\real \subseteq L^2(\real)\).
    \begin{enumerate}[label=(\alph*), noitemsep]
        \item Show that \(\real\) and \(\real^\perp \coloneq \set{[X] \in L^2 :
        \E{X}=0}\) are orthogonal subspaces of \(L^2\) with \(L^2 = \real
        \oplus \real^\perp\).
        \begin{solution}
            Let \(r\in \real\) and \([X]\in \real^\perp\), then
            \[
                \scp{[r], [X]}_{L^2} = \E{rX} = r\E{X} = 0.
            \]
            Consequently \(\real \perp \real^\perp\).  Moreover for any \([X]
            \in L^2\), we have
            \[
                X = \underbrace{\E{X}}_{\in \real} + \underbrace{(X-\E{X})}_{\in \real^\perp}.
            \]
            Thus \(L^2 = \real \oplus \real^\perp\).
        \end{solution}

        \item Show that the orthogonal projection onto \(\real\) is given by
        \[
            \eSym \colon \begin{cases}
                L^2 \to \real,
                \\
                [X] \mapsto \E{X}.
            \end{cases}
        \]
        \begin{solution}
            Since \(\E{\E{X}} = \E{X}\) we have \(\eSym^2 = \eSym\)
            and therefore the expectation is a projection as it is also a
            linear operator. It is an orthogonal projection since
            \[
                \scp{\E{X}, [Y]}_{L^2} = \E{\E{X}Y} = \E{X}\E{Y} = \scp{[X], \E{Y}}_{L^2}.
                \qedhere
            \]
        \end{solution}

        \item Let \(\pi \colon L^2 \to \real^\perp\) be the orthogonal projection
        onto \(\real^\perp\), that is \(\pi = 1 - \eSym\). Show that
        \[
            \cov{X}{Y}
            = \scp{\pi[X], Y}_{L^2}
            = \scp{X, \pi[Y]}_{L^2}
            = \scp{\pi[X], \pi[Y]}_{L^2}
        \]
        In particular \(\covSym\) coincides with the inner product on \(\real^\perp\)!
        \begin{solution}
            By definition of the covariance
            \[
                \cov{X}{Y}
                = \E{(X-\E{X})(Y-\E{Y})}
                = \E{\pi[X]\pi[Y]} = \scp{\pi[X], \pi[Y]}_{L^2}.
            \]
            The other equalities essentially follow from the fact that this is an
            orthogonal projection:
            \[
                \scp{\pi[X], Y}_{L^2}
                = \scp{\pi[X], \E{Y} + \pi[Y]}_{L^2}
                = \underbrace{\scp{\pi[X], \E{Y}}_{L^2}}_{=0} + \scp{\pi[X], \pi[Y]}_{L^2}
            \]
            This can also be seen by direct calculation:
            \[
                \scp{\pi[X], \E{Y}}_{L^2} = \E{(X-\E{X})\E{Y}} = \E{X -\E{X}}\E{Y} = 0.
                \qedhere
            \]
        \end{solution}

        \item Show that the expectation is the closest real number to \(X\in
        \cL^2\) in the \(L^2\) distance, that is
        \[
            \E{X} = \argmin_{c\in \real} \E{(X - c)^2} = \argmin_{c\in \real} \norm{X - c}_{L^2}.
        \]
        \begin{solution}
            This is a property of the orthogonal projection. Specifically,
            \[
                \norm{X - c}_{L^2}^2
                = \norm{\pi[X] + \E{X} - c}_{L^2}^2
                = \norm{\pi[X]}_{L^2}^2 + \norm{c-\E{X}}_{L^2}^2,
            \]
            as the mixed terms disappear, indeed:
            \[
                \scp{\pi[X], c-\E{X}}_{L^2} = \E{X - \E{X}}(c-\E{X}) = 0.
            \]
            Since \(\norm{\pi[X]}_{L^2}^2\) does not depend on \(c\) and the
            square is a monotone function, we consequently have
            \[
                \argmin_{c\in \real} \norm{X - c}_{L^2}
                = \argmin_{c\in \real} \norm{c-\E{X}}_{L^2}
                = \argmin_{c\in \real} \abs{c-\E{X}}
                = \E{X}.
                \qedhere
            \]
        \end{solution}
    \end{enumerate}
    \begin{remark}[Generalization to Hilbert spaces]
        The same arguments work if we replace \(\real\) by any Hilbert space
        \((\hilbert, \scp{\cdot, \cdot})\).
    \end{remark}
\end{exercise}

\begin{exercise}[Median -- Projection in \(L^1\) \Coffeecup\Coffeecup]
	We define the median of \(X\) to be the set
    \[
        \mSym[X] \coloneq \set[\Big]{m \in \real : \prob{X \ge m}\ge \tfrac12, \prob{X\le m} \ge \tfrac12}.
    \]
	\begin{enumerate}[label=(\alph*), noitemsep]
        \item Show that if \(a,b\in \mSym[X]\), then \((a,b) \subseteq \mSym[X]\).
        That is \(\mSym[X]\) is an interval!

        \item The next goal is to show that the median interval is closed
        and characterize the upper and lower bound. Prove for \(F(x) = \prob{X\le x}\)
        \begin{alignat*}{2}
            \max \mSym[X] &= \inf \set{x \in \real : F(x) > \tfrac12}
            &&\eqcolon \overline{M},
            \\
            \min \mSym[X] &= \inf \set{x \in \real : F(x) \ge \tfrac12} \overset{\text{def.}}= F^{\leftarrow}(\tfrac12)
            &&\eqcolon \underline{M}.
        \end{alignat*}
        This has two components: \(\overline{M}\) and \(\underline{M}\) are the boundaries of \(\mSym[X]\)
        and they are contained in \(\mSym[X]\).

        \textbf{Hint:} Use \(\prob{X < x} = \prob{\bigcup_{k\in \nat}\set{X \le x - \tfrac1k}}\) and the continuity of measure.
        \begin{solution}
            Since the distribution function \(F\) is right-continuous we have
            \[
                \prob{X \le \overline{M}} = F(\overline{M}) = \lim_{x \downarrow \overline{M}} F(x) \ge \tfrac12
            \]
            as all \(x \ge \overline{M}\) have the property \(F(x) > \tfrac12\) by
            definition of \(\overline{M}\). For the other direction we use the
            continuity of measure:
            \[\begin{aligned}
                \prob{X \ge \overline{M}}
                &= 1 - \prob{X < \overline{M}}
                \\
                &= 1- \prob[\Big]{\bigcup_{k\in \nat} \set{X \le \overline{M} - \tfrac1k}}
                \\
                &= 1 - \lim_{k\to \infty} F(\overline{M}-\tfrac1k)
                \\
                &\ge 1- \tfrac12 = \tfrac12
            \end{aligned}
            \]
            as any \(x< \overline{M}\) has the property \(F(x) \le \tfrac12\).
            Consequently, \(\overline{M} \in \mSym[X]\).
            As we have not used the strict inequalities due to the limits, the
            same arguments work for the minimum \(\underline{M} = \inf\set{x\in
            \real F(x) \ge \tfrac12}\). We only replace \(F(x) > \tfrac12\) by
            \(F(x) \ge \tfrac12\) for all \(x \ge \underline{M}\) for the first
            part and use \(F(x) < \tfrac12\) for all \(x < \underline{M}\) for
            the second. Now by definition of \(\overline{M}\) we have for any
            \(m > \overline{M}\) that \(F(m) > \tfrac12\) and therefore
            \[
                \prob{X \ge m} = 1 - F(m) < \tfrac12 \qquad \forall m > \overline{M}.
            \]
            This implies \(\overline{M} = \max \mSym[X]\).
            Similarly, we have for any \(m < \underline{M}\) that \(F(m) < \tfrac12\)
            and therefore \(\underline{M} = \min \mSym[X]\).
        \end{solution}

        \item\label{it: characterization of distribution function} Prove that
        \[
            \begin{cases}
                F(x) > \tfrac12 & x > \overline{M} 
                \\
                F(x) = \tfrac12 & x \in (\underline{M}, \overline{M})
                \\
                F(x) < \tfrac12  & x < \underline{M}
            \end{cases}
        \]
        In particular, if the distribution function \(F\) jumps over
        \(\tfrac12\) such that there exists no \(x\) with \(F(x) = \frac12\),
        then \(\underline{M} = \overline{M}\). For the case
        \(\underline{M} \neq \overline{M}\) prove \(F(\underline{M})= \tfrac12\).
        \begin{solution}
            \(F(x) > \tfrac12\) for \(x>\overline{M}\) follows from the definition
            of \(\overline{M}\). Similarly \(F(x) < \tfrac12\) for \(x<\underline{M}\)
            follows from the definition of \(\underline{M}\). Now for any
            \(x \in (\underline{M}, \overline{M})\) we have \(x> \underline{M}\) and
            consequently \(F(x) \ge \tfrac12\). And since \(F(x) > \tfrac12\) would imply
            \(x > \overline{M}\) by definition of \(\overline{M}\) we must conclude
            \(F(x) = \tfrac12\). Finally if \(\underline{M} \neq \overline{M}\)
            then \(F(\underline{M}) = \tfrac12\) follows from right-continuity of \(F\).
        \end{solution}

		\item Assuming \(X\in \cL^1\) prove that the median minimizes the \(L^1\) error, i.e.
        \[
            \mSym[X] =\arg\min_{m} \E{\abs{X-m}}.
        \]
		% If you wish to prove the general
		% case, note that \(f(m)=2F_X(m)-1\) is right-continuous and monotonously
		% increasing. Monotonicity also implies it can only have a countable number
		% of discontinuities.

 		\textbf{Hint:}
            Use \(\E{X}=\int_0^\infty \prob{X\ge t}dt = \int_0^\infty \prob{X>t}dt\)
            for \(X\ge 0\) to prove
            \[
                \E{\abs{X-m_2}} - \E{\abs{X-m_1}}
                = \int_{m_1}^{m_2} \underbrace{\prob{X\le x} - \prob{X>x}}_{=2F(x)-1}dx
            \]
		\begin{solution}
			We have
			\begin{align*}
				\E{\abs{X-m}_+}
				&= \int_0^\infty \underbrace{\prob{\abs{X-m}_+ > t}}_{
					\begin{aligned}
						&=\prob{X-m>t}\\
						&=\prob{X> t-m}
					\end{aligned}
				} dt\\
				&= \int_m^\infty \prob{X>s}ds
			\end{align*}
			And similarly
			\begin{align*}
				\E{(X-m)_-}
				&= \int_0^\infty \underbrace{\prob{(X-m)_- \ge t}}_{
					\begin{aligned}
						&=\prob{m-X\ge t} \quad \mathrlap{(\forall t>0)}\\
						&=\prob{m-t\ge X}
					\end{aligned}
				} dt\\
				&= \int_{-\infty}^m \prob{X\le s}ds
			\end{align*}
			Put together we have
            \[
				\E{\abs{X-m}}
				= \int_{-\infty}^m \prob{X\le s}ds + \int_m^\infty \prob{X>s}ds.
            \]
            Consequently, we have for \(m_1 < m_2\)
            \[\begin{aligned}
                \E{\abs{X-m_2}} - \E{\abs{X-m_1}}
                &= \int_{m_1}^{m_2} \prob{X\le s}ds - \int_{m_1}^{m_2} \prob{X>s}ds
                \\
                &= \int_{m_1}^{m_2} \underbrace{\prob{X\le s} - (1-\prob{X\le s})}_{=2F(x) -1}ds
            \end{aligned}
            \]
            Using \ref{it: characterization of distribution function} we deduce
            \begin{alignat*}{3}
                \E{\abs{X-m}} - \E{\abs{x-\overline{M}}} &> 0
                & m &> \overline{M} 
                \\
                \E{\abs{X-m_2}} - \E{\abs{x-m_1}} &= 0
                \qquad & m_1, m_2 &\in (\underline{M}, \overline{M})
                \\
                \E{\abs{x-\underline{M}}} - \E{\abs{X-m}} &< 0
                & m &< \underline{M}
            \end{alignat*}
            Reordering implies the claim. \qedhere
		\end{solution}
	\end{enumerate}
\end{exercise}

\begin{remark}[Median is more stable]
    Observe that both the median and expectation are defined in greater
    generality than their interpretation as projections permits
    \begin{center}
       \def\arraystretch{1.2} 
    \begin{tabular}{l|ll}
        & expectation & median 
        \\
        \hline
        defined & in \(L^1\) & always
        \\
        minimizes & \(L^2\)-norm & \(L^1\)-norm
    \end{tabular}
    \end{center}
    However the median is more stable and works for heavy tailed
    distributions that do not have a well defined expectation -- it is
    more robust to outliers.
\end{remark}

\begin{remark}[Median in higher dimension]
    The median is much more difficult to work with than the expectation
    as it is set-valued in general and more difficult to compute.
    In higher dimension it is even more difficult to define as there is no
    natural ordering and consequently no distribution function or quantile
    function. Finding a good generalization is still an active area of
    research! A recent approach \citep{hallinDistributionQuantileFunctions2021}
    exploits the idea that the distribution function is a monotonous function.
    More specifically, the shifted function \(F(x) - \frac12\) is decreasing
    before the median and increasing after the median. Integrating this function
    therefore gives a convex function that is minimized at the median.
    While a monotonous function is only defined on ordered sets, convexity
    generalizes to higher dimensions. This idea is used in
    \citep{hallinDistributionQuantileFunctions2021} to define higher dimensional
    quantile functions as the unique gradient of a convex function
    that satisfies Property \ref{it: inverse transform sampling} of Exercise \ref{ex:
    stochastic simulation}.
\end{remark}